\documentclass[UTF8]{article}
\usepackage[a4paper,margin=1in]{geometry}
\usepackage{amsmath}
\usepackage{amssymb}
\usepackage{graphicx}
\usepackage{booktabs}
\usepackage{longtable}
\usepackage{array}
\usepackage{hyperref}
\usepackage{xeCJK}

\begin{document}

% Page 1
\noindent 第五章 \quad 图像复原 \hfill 河海大学

\begin{center}
    \Huge \textbf{第五章 \quad 图像复原}
\end{center}

\vspace{1cm}

\noindent \Large \textbf{主要内容：}

\begin{itemize}
    \item \Large \textbf{5. 1. \quad 概述}
    \item \Large \textbf{5. 2. \quad 图像降质的数学模型}
    \item \Large \textbf{5. 3. \quad 无约束图像复原}
    \item \Large \textbf{5. 4. \quad 约束图像复原}
    \item \Large \textbf{5. 5. \quad 典型的几何失真校正}
    \item \Large \textbf{5. 6. \quad 图像几何空间变换}
\end{itemize}

\vfill
\noindent 数字图像处理 \hfill $<$ \quad 1 \quad $>$

% Page 2
5.1 概述 \hfill 河海大学

\section*{5.1. 图像复原概述}

图像降质

\begin{center}
    \includegraphics[width=0.8\textwidth]{image_process_diagram.png} % 假设图片文件名为 image_process_diagram.png
    \\
    如何实现恢复？
\end{center}

\noindent 脉冲噪声（椒盐噪声） \hfill 复原后图像 \\
污染后的图像

\noindent 数字图像处理 \hfill < \quad 2 \quad >

% Page 3
5.1 概述 \hfill 河海大学

\section*{图像降质}

\begin{center}
    % 此处为示意图占位
    \text{运动形成的模糊} \quad \xrightarrow{\text{如何实现恢复？}} \quad \text{复原后图像}
\end{center}

数字图像处理 \hfill < \quad 3 \quad >

% Page 4
5.1 概述

图像降质

\begin{center}
    \includegraphics[width=0.8\textwidth]{image_process_diagram.png} % 假设图片文件名为 image_process_diagram.png
\end{center}

\begin{minipage}{0.45\textwidth}
    \centering
    离焦形成的模糊
\end{minipage}
\hfill
\begin{minipage}{0.45\textwidth}
    \centering
    复原后图像
\end{minipage}

数字图像处理 \hfill < \quad 4 \quad >

% Page 5
5.1 概述

\section*{图像降质}

\begin{itemize}
    \item 噪声等形成的降质
    \item 运动引起的降质
    \item 亚采样引起的降质
\end{itemize}

\begin{center}
    \fbox{\parbox{0.25\textwidth}{\centering \color{red} 图像的降质 \\ 或者退化}}
\end{center}

\begin{enumerate}
    \item 成像系统的像差、畸变、有限带宽等造成的图像失真；
    \item 涉嫌辐射、大气测流等造成的照片畸变；
    \item 携带遥感仪器的飞机或卫星运动的不稳定，以及地球自传等因素引起的照片几何失真；
    \item 模拟图像在数字化的过程中，由于会损失掉部分细节，因而造成图像质量下降；
    \item 拍摄时，相机与景物之间的相对运动产生的运动模糊；
    \item 镜头聚焦不准产生的散焦模糊；
    \item 底片感光、图像显示时造成记录显示失真；
    \item 成像系统中存在的噪声干扰。
\end{enumerate}

数字图像处理 \hfill 5
\end{document}
